\section*{Your turn}

Chapter~\ref{ch:why-one-graph-is-never-enough}'s exercise needed no code. This
one does. Clone the companion repository and check out the tag for this chapter:

\begin{minted}{text}
git clone https://github.com/Giang-Dang/mosaic-graph
cd mosaic-graph
git checkout ch02
pwsh scripts/verify.ps1
\end{minted}

The script builds the solution, checks the committed schema against a freshly
exported one, starts the service, and runs the Postman collection. If it prints
\code{PASS} you have the same Mosaic I have. Appendix~\ref{app:toolchain} covers
the setup in more detail and gives you a tour of the repository.

\begin{enumerate}
  \item \textbf{Prove the styles are interchangeable.} Run all three projects
  under \code{samples/three-approaches} and export a schema from each. Diff
  them, or compare their hashes. Then break one deliberately: rename a field in
  the code-first descriptor and watch which of the three schemas moves. The
  point of the exercise is the direction of the arrow, from code to schema, and
  how little of it a client can see.

  \item \textbf{Add a field, and watch the diff.} Give \code{Product} a
  \code{weightInGrams}. Put it wherever you think it belongs, then export the
  schema and look at the diff. Where did the field land in the type, and does
  that match where you would have written it by hand?

  \item \textbf{Break the schema on purpose.} Give two resolvers on the same
  type the same GraphQL name. Start the service. Note where the failure happens
  and what it costs you: with eager initialisation this is a process that will
  not boot, not a query that returns an error. Now imagine that container going
  out in a rolling deploy.

  \item \textbf{Move the number.} Section~\ref{sec:ch02-still-wrong} counts 146
  lookups for the catalog-with-reviews query. Work out, without running it, what
  the count should be for a query that asks for every product's title and price
  and nothing else. Then run it and check. If you were wrong, the interesting
  question is which resolver you forgot.

  \item \textbf{Make a lookup cost something.} In \code{appsettings.json}, set
  the lookup delay under \code{Mosaic:Data} to five milliseconds and re-run the
  catalog-with-reviews query. The service does exactly what it did before and
  the schema is unchanged. Write down the new response time before you look at
  it, then look.

  \item \textbf{Lose introspection.} Start the service with
  \code{ASPNETCORE\_ENVIRONMENT=Production} and try to refresh the schema in
  Postman. Read the error. This is the failure mode most likely to greet you the
  first time you point a tool at a deployed graph, and it is much less alarming
  when you have seen it on purpose.
\end{enumerate}

Exercise~5 is the one to sit with. Five milliseconds is a plausible round trip
to a database on the same network, and it is a number nobody would flag in
review. Multiply it by the count from exercise~4 and you have the reason
chapter~\ref{ch:data-without-the-n-plus-1} exists.
